\documentclass[12pt]{article}
\usepackage[utf8]{inputenc}
\usepackage[T1]{fontenc}
\usepackage{hyperref}
\title{Zorop Language Resources for K-12 Curriculum Development}
\author{Compiled by Claude Sonnet 5.5 (claude-sonnet-5-5) for Andr\'as Kornai}
\date{}
\begin{document}
\maketitle

\section{Introduction}
Zorop (also Yafi, Jafi, Warlef or Yaffi; ISO 639-3 wfg) is an East Pauwasi language. It is spoken in Warlef village, Senggi District, Keerom Regency, Indonesia. Speaker counts differ between sources: about 230 (2005) and about 300. Wikitongues lists it as EGIDS 7 (shifting). Very little material exists.

\section{Primary Educational Resources}
No primers, textbooks or school programmes were found in search results.

\section{Linguistic and Reference Materials}
\begin{itemize}
\item Wikipedia article on the Yafi language, with a short vocabulary table. \url{https://en.wikipedia.org/wiki/Yafi_language}
\item Wikipedia articles on the Pauwasi and East Pauwasi languages. \url{https://en.wikipedia.org/wiki/Pauwasi_languages} and \url{https://en.wikipedia.org/wiki/East_Pauwasi_languages}
\item University of Zurich New Guinea World page on East Pauwasi River languages. \url{https://newguineaworld.linguistik.uzh.ch/_export/raw/families/pauwasi-river/east-pauwasi-river/start}
\item Indonesian Wikipedia article on Zorop. \url{https://id.wikipedia.com/wiki/Bahasa_Zorop}
\end{itemize}

\section{Religious and Community Resources}
\begin{itemize}
\item Joshua Project profile of the Yafi people (demographic information only). \url{https://joshuaproject.net/people_groups/15937}
\end{itemize}
No scripture or community organisation was found.

\section{Digital and Audio Resources}
\begin{itemize}
\item Wikitongues page for wfg. No recordings were confirmed in the search results. \url{https://wikitongues.org/languages/wfg/}
\end{itemize}

\section{Curriculum Development Context}
The language is endangered and used as a first language by adults only. No written materials for teaching were found. Any curriculum would have to start from new documentation done with the community.

\end{document}